\section{The conclusions}\label{sec:result}

\noindent\textbf{C1. No worldview can be established as the true account.}\phantomsection\label{item:C1}
(From D1, D3, P-M(b) and AA-P1.)
Let W be a worldview.
\begin{itemize}
  \item W goes beyond the evidence (D1).
  \item By P-M(b), a rival R consistent with the evidence and the truths of reason makes W false.
  \item By D3, no argument establishes W, since W is false in R.
  \item This holds at every time, because evidence stays finite (AA-P1).
\end{itemize}
No premise about premises transmitting establishment is needed, since D3 concerns only the
accounts consistent with the evidence.

Stated for the worldviews at issue:

\begin{quote}
No argument from experience or reason establishes that any worldview, whether theism,
naturalism or the skeptical hypothesis, is the true account of what is. Any of them may be true,
but no one can establish that it is.
\end{quote}

This is agnosticism in Huxley's sense, ``not a creed, but a method'': ``do not pretend that
conclusions are certain which are not demonstrated or demonstrable''
\citep[pp.~245--246]{huxley1889}. The claim concerns what can be established, not what it is
reasonable to believe.

\medskip\noindent\textbf{C2. Worldviews that agree on all possible evidence are never
separated.}\phantomsection\label{item:C2}
(From P-M(a) and AA-R2.)
By P-M(a), two worldviews that agree on everything that could ever be experienced are accounts
that agree on all possible evidence. By AA-R2, no further evidence and no ideal judge ever
separates them. The accounts of \emph{Algorithmic Agnosticism} do not speak of God or nature, so applying AA-R2
to worldviews is this paper's extension, made through P-M(a). The hope that future or final
evidence would settle the matter therefore fails for such pairs.

Hick's eschatological verification is a case of that hope. Theism, Hick argued, is verifiable in
principle by experience after death, such as the experience of God's purpose for human beings
fulfilled in communion with God as revealed in Christ, so theism and naturalism are not empirically equivalent
\citep{hick1960}. C2 does not decide whether he was right; it says what follows on each answer. If a
non-theistic account could accommodate every such experience, it and theism agree on all possible
evidence, and C2 applies: the experience, should it come, would not separate them. If not, the two
differ on possible evidence, and C2 does not apply to them. The experience, when it came, would
still be finite evidence with rivals consistent with it, and C1 answers any claim that it
established theism.

\medskip\noindent\textbf{C3. What selects among the worldviews the evidence leaves open is a
choice.}\phantomsection\label{item:C3}
(From P-M(a) and AA-R3.)
By P-M(a), worldviews are accounts, so by AA-R3 a prior probability, simplicity, or any other
criterion that ranks the rivals is an imposed constraint, and the evidence does not determine it.
Probabilistic and parsimony arguments therefore establish nothing beyond the constraint they adopt.

Each decisive element in \cref{sec:arguments} is thus answered: an account or a judgment by C1,
and an imposed constraint, a prior included, by C3.

The conclusions are conditional on P-M, the premise AA-P1 of \emph{Algorithmic Agnosticism}, and logic. D1 and D3
are definitions, and its results are proved there. The conclusions concern
arguments and evidence, not what ultimately is. P-M is not established by D3's standard either;
it is owned as a premise, and the inference from it is logic.
